\documentclass[a4paper, 11pt]{article}

% --- PACOTES ---
\usepackage[utf8]{inputenc}
\usepackage[brazil]{babel}
\usepackage{geometry}
\usepackage{hyperref}
\usepackage{longtable}
\usepackage{booktabs}
\usepackage{array}
\usepackage{ragged2e}
\usepackage{graphicx}
\usepackage{booktabs} 
\usepackage{threeparttable}
\usepackage{graphicx} 
\usepackage{siunitx}
\sisetup{round-mode=none}
\usepackage{adjustbox}


\usepackage[T1]{fontenc}
\usepackage{csquotes}

\usepackage[backend=biber, style=authoryear, citestyle=authoryear]{biblatex}
\DeclareDelimFormat{nameyeardelim}{\addcomma\space}
\setlength{\bibitemsep}{1em}
\addbibresource{bibliografia.bib}

% --- CONFIGURAÇÕES DO DOCUMENTO ---
\geometry{a4paper, margin=1in}
\hypersetup{
    colorlinks=true,   
    citecolor=black,   
    linkcolor=black,
    urlcolor=blue 
}
\setlength{\parindent}{1.25cm}
\setlength{\parskip}{0.5em}

\begin{document}
\begin{titlepage}

    \noindent\includegraphics[width=0.22\textwidth]{ufsc.png}\par
    \vspace{0.5cm}
    
    \centering
    \textbf{Universidade Federal de Santa Catarina - UFSC}\\
    Departamento de Informática e Estatística (INE)\\
    Disciplina INE5429 - Segurança em Computação\\[0.5cm]
    Caio Prá Silva (21203773)\\
    
    \vfill
    
    \centering
    {\bfseries\Large
    Relatório Trabalho Individual - Números primos\par}
    
    \vfill 
    
    \centering
    \textbf{Professores:} Prof. Jean Everson Martina e Prof$^{a}$. Thaís Bardini Idalino\\[0.3cm]
    
    \centering
    \textbf{Florianópolis, SC}\\
    \textbf{2025}
    
\end{titlepage}

% --- SEÇÃO 1 ---
\section{Introdução}

\subsection{Motivação e Fundamentação Teórica}

A segurança da informação na era digital é amplamente sustentada por protocolos criptográficos robustos, muitos dos quais dependem fundamentalmente de conceitos da teoria dos números. No centro de grande parte da criptografia de chave pública moderna, como o sistema RSA (Rivest-Shamir-Adleman), reside a utilização de números primos de grande magnitude. \parencite{rsa-aes} A segurança destes sistemas deriva da assimetria computacional entre a facilidade de multiplicar dois números primos grandes e a extrema dificuldade de fatorar o seu produto.  Para aplicações práticas, como a assinatura digital de documentos eletrônicos no Brasil, são exigidas chaves criptográficas geradas a partir de números primos com bits ou mais, uma dimensão que impõe desafios computacionais significativos. \parencite{iti_doc_icp_01_01}

O processo de obtenção de um número primo de grande porte para uso criptográfico é tipicamente decomposto em duas tarefas sequenciais e interdependentes. A primeira consiste na geração de um número candidato, que deve ser um inteiro ímpar de magnitude apropriada e possuir características de aleatoriedade para evitar ataques preditivos. A segunda tarefa é a verificação da primalidade deste candidato, um processo que deve ser tanto eficiente quanto confiável, dado que um teste incorreto poderia comprometer toda a segurança do sistema criptográfico subsequente. Esta decomposição do problema evidencia a necessidade de se dispor de algoritmos eficientes e robustos para ambas as etapas: a geração de números pseudo-aleatórios e o teste de primalidade.

O processo de obtenção de um número primo de grande porte para uso criptográfico é tipicamente decomposto em duas tarefas sequenciais e interdependentes. A primeira consiste na geração de um número candidato, que deve ser um inteiro ímpar de magnitude apropriada e possuir características de aleatoriedade para evitar ataques preditivos. A segunda tarefa é a verificação da primalidade deste candidato, um processo que deve ser tanto eficiente quanto confiável, dado que um teste incorreto poderia comprometer toda a segurança do sistema criptográfico subsequente. Esta decomposição do problema evidencia a necessidade de se dispor de algoritmos eficientes e robustos para ambas as etapas: a geração de números pseudo-aleatórios e o teste de primalidade.

\subsection{Objetivos}
Este relatório tem como objetivo principal conduzir uma análise teórica e empírica de algoritmos selecionados para cada uma dessas tarefas, com foco particular na sua aplicação a inteiros de precisão arbitrária, capazes de atingir os 4096 bits exigidos pelo escopo do trabalho. Para a geração de números pseudo-aleatórios, serão implementados e comparados o Gerador Linear Congruente (LCG) \parencite{lehmer1951}, um método clássico baseado em aritmética modular, e o gerador Xorshift \parencite{xorshift-def}, uma abordagem moderna otimizada para operações a nível de bit. Para a verificação de primalidade, o trabalho implementará o Teste de Miller-Rabin \parencite{miller1976, rabin1980}, um padrão da indústria e requisito obrigatório, e o comparará com Teste de Fermat \parencite{solovay1977}, um precursor historicamente importante cujas limitações teóricas oferecem um contraponto valioso para a análise. A análise comparativa abrangerá o desempenho computacional (tempo de execução), a complexidade assintótica e a robustez dos algoritmos, avaliando a qualidade estatística dos geradores e a correção dos testes de primalidade.

\subsection{Considerações sobre Aritmética de Múltipla Precisão}

O requisito de gerar números de até 4096 bits  excede em muito a capacidade dos tipos de dados inteiros nativos da maioria das linguagens de programação, como int (32 bits) ou long long (64 bits) em C/C++. A manipulação de números dessa magnitude requer o uso de aritmética de precisão arbitrária, comumente conhecida como "\textit{bignum}" \parencite{oshiro2020how_python_bignum}. 

Para este trabalho, a linguagem de programação Python foi escolhida devido ao seu suporte nativo e transparente a inteiros de precisão arbitrária. Em Python, o tipo int pode representar números de qualquer tamanho, limitado apenas pela memória disponível no sistema. Essa característica abstrai a complexidade do gerenciamento de memória e da implementação de operações aritméticas para números grandes, permitindo que o foco permaneça na lógica dos algoritmos de PRNG e primalidade. É importante notar que, embora conveniente, a implementação de \textit{bignum} nativa do Python pode não ser a mais performática. Para aplicações de alta performance em ambientes de produção, bibliotecas especializadas escritas em C/C++, como a GNU Multiple Precision Arithmetic Library (GMP) \parencite{gmp} , a Crypto++  ou a \textit{s2n-bignum} da \textit{AWS} \parencite{awslabs_s2n_bignum} , ofereceriam um desempenho superior. O Python pode interagir com bibliotecas através de wrappers como o gmpy2 , que combina a facilidade do Python com velocidade do C. No entanto, para o escopo deste trabalho, o desempenho do Python nativo é suficiente e a simplicidade que ele oferece é uma vantagem significativa.

% --- SEÇÃO 2 ---
\section{Geração de Números Pseudo-Aleatórios (\textit{PRNGs})}

A geração de números pseudo-aleatórios (\textit{Pseudo-Random Number Generators} - \textit{PRNG}) é o primeiro passo crítico no processo de encontrar um primo grande para fins criptográficos. A qualidade do \textit{PRNG} é fundamental; uma sequência previsível pode permitir que um adversário restrinja o espaço de busca de chaves, comprometendo a segurança do sistema. Nesta seção, dois algoritmos com de projeto distintas são explorados: o Gerador Linear Congruencial (\textit{linear congruential generator} - \textit{LCG}), que se baseia em aritmética modular, e o Xorshift, que utiliza operações \textit{bitwise}.

\subsection{Gerador Linear Congruencial - \textit{LCG}}
\label{sub-lcg}

O Gerador Linear Congruencial é um dos algoritmos de PRNG mais antigos e estudados na computação. Sua popularidade histórica deve-se à sua simplicidade de implementação e à sua velocidade em hardware que suporta nativamente a aritmética modular \parencite{wikipedia_lcg}. A operação do LCG é definida pela seguinte relação de recorrência:
\begin{equation}
    X_{n+1} = (a \cdot X_n +c) \quad (mod \; m)
\end{equation}
\begin{flushleft}
    onde:
    \begin{itemize}
        \item \textbf{$X_n$} \textit{é o n-ésimo} número da sequência pseudo-aleatória;
        \item \textbf{$X_0$} é a "semente" (\textit{seed}), o valor inicial que define toda a sequência subsequente;
        \item \textbf{$a$} é o "multiplicador" $\,(0 < a< m)$;
        \item \textbf{$c$} é o "incremento" $\,(0 < c< m)$;
        \item \textbf{$m$} é o "módulo" $\,(m>0)$, que define o intervalo dos números gerados, $[0, m-1]$.
    \end{itemize}
\end{flushleft}

A qualidade de um gerador \textit{LCG} é muito sensível à escolha de seus parâmetros. Para que o gerador atinja seu período máximo, é necessário que:

\begin{enumerate}
    \item c e m devem ser coprimos
    \item $a-1$ deve ser múltiplo de todo fator primo de m
    \item $a-1$ deve ser múltiplo de 4 se m for múltiplo de 4
\end{enumerate}

A história da computação está repleta de exemplos onde a escolha inadequada de parâmetros levou a geradores com falhas estatísticas. Um caso notório é o \textit{RANDU} \parencite{ibm1968randu}, um LCG amplamente utilizado na década de 1970, cujos parâmetros resultavam em uma forte correlação entre valores sucessivos, invalidando muitas simulações científicas que o utilizaram. \parencite{marsaglia1968}

\subsection{Gerador Xorshift}

Os geradores Xorshift, introduzidos por \cite{JSSv008i14}, representam uma classe de \textit{PRNGs} que se destacam pela sua simplicidade e velocidade excepcional em implementações de software. Eles pertencem à família mais ampla de registradores de deslocamento com linear (LFSRs), mas são projetados para serem particularmente eficientes em arquiteturas de computadores modernas, que executam operações \textit{bitwise} (a nível de bit) em poucos ciclos de \textit{clock}. \parencite{wikipedia_xorshift}

O mecanismo central do Xorshift consiste em gerar o próximo número de uma sequência aplicando repetidamente a operação "ou exclusivo" (XOR, denotada por $\oplus$) a um estado interno com versão deslocada de si mesmo. Uma implementação típica de 32 bits, por exemplo, pode seguir a forma:

$$x = x \,\oplus\, (x << a);$$
$$x = x \,\oplus\, (x >> b);$$
$$x = x \,\oplus\, (x << c);$$

\begin{flushleft}
    onde $x$ é o estado do gerador (inteiro de 32 bits), $<<$ e $>>$ são as operações de deslocamento de bits para a esquerda e para a direita, respectivamente, e $a$, $b$ e $c$ são constantes de deslocamento escolhidas. A escolha de tais parâmetros é crucial para garantir maior periodicidade para o gerador, bem como garantir a ele boas propriedades estatísticas.
\end{flushleft}

Embora os geradores Xorshift puros sejam extremamente rápidos, eles podem falhar em baterias de testes estatísticos mais rigorosas, como o BigCrush. Para mitigar essa fraqueza, foram propostas variações que combinam a saída do Xorshift com função não-linear. As variantes mais comuns incluem xorshift*, que aplica uma multiplicação modular, e xorshift+, que utiliza uma adição, melhorando significativamente a qualidade estatística da sequência gerada. Para os propósitos deste trabalho, a implementação da versão pura do Xorshift é suficiente e oferece um excelente ponto de comparação com LCG, destacando o \textit{trade-off} entre complexidade aritmética e operações bitwise. \parencite{vigna_xorshift_prng}

\subsection{Implementação e Metodologia Experimental}

A implementação e a experimentação foram projetadas para avaliar os algoritmos no contexto de números inteiros de grande porte, conforme exigido pelas aplicações criptográficas modernas.

\subsubsection{Configuração Experimental}
\label{sec:config-exp}

Os experimentos foram conduzidos em um sistema com seguintes especificações:
\begin{itemize}
    \item \textbf{CPU:} Intel Core i5-11400H @ 2.70GHz
    \item \textbf{RAM:} 16 GB
    \item \textbf{SO:} Ubuntu 22.04
    \item \textbf{Versão Python:} 3.10.12
\end{itemize}

Para medir o desempenho, foi utilizada a biblioteca \textit{\textbf{time}} do Python, especificamente com método perf\_counter \parencite{python_time_perf_counter}, para medições de tempo mais precisas durante as múltiplas execuções de código, minimizando a influência de flutuações do sistema. Conforme sugerido na especificação do trabalho, usou-se tamanho de bits: 40, 56, 80, 128, 168, 224, 256, 512, 1024, 2048, 4096. Foram gerados 16 valores para cada quantidade de dígitos binários, com mínimo, máximo e médio de geração, bem como desvio padrão e os tempos individuais.

Especificamente para o algoritmo LCG, foram usados os parâmetros baseados na implementação ANSI C \parencite{saucier2000computer}:
\begin{itemize}
    \item m = $2^{32}$
    \item a = $1103515245$
    \item c = $12345$
\end{itemize}

\subsubsection{Experimentos}

Considerando os tamanhos de bits mencionados e demais especificações, usando o código disponível no GitHub (Apêndice \ref{app-a}), obteve-se os resultados sumarizados na Tabela \ref{tab:tempos1} e Tabela \ref{tab:tempos2}.

\begin{table}[h!] \centering \begin{threeparttable} \caption{Resultados de desempenho do LCG} \label{tab:tempos1} \begin{tabular}{|c|c|c|c|c|} \hline \textbf{N bits} & \textbf{Tempo médio (ms)} & \textbf{Tempo mín (ms)} & \textbf{Tempo máx (ms)} & \textbf{Desvio padrão (ms)} \\ \hline 40 & 0.001994 & 0.001058 & 0.009855 & 0.002345 \\ \hline 56 & 0.001103 & 0.001023 & 0.001547 & 0.000126 \\ \hline 80 & 0.001389 & 0.001319 & 0.001821 & 0.000123 \\ \hline 128 & 0.001720 & 0.001642 & 0.002125 & 0.000112 \\ \hline 168 & 0.002301 & 0.002205 & 0.002741 & 0.000125 \\ \hline 224 & 0.002625 & 0.002547 & 0.003072 & 0.000123 \\ \hline 256 & 0.002914 & 0.002845 & 0.003270 & 0.000101 \\ \hline 512 & 0.005495 & 0.005277 & 0.006719 & 0.000345 \\ \hline 1024 & 0.010003 & 0.009270 & 0.012404 & 0.000877 \\ \hline 2048 & 0.019950 & 0.018419 & 0.034388 & 0.004017 \\ \hline 4096 & 0.042365 & 0.039881 & 0.049948 & 0.003641 \\ \hline \end{tabular} \begin{tablenotes} \small \item Os valores representam a média, mínimo, máximo e desvio padrão de 16 execuções independentes para cada tamanho de chave. \end{tablenotes} \end{threeparttable} \end{table} \begin{table}[h!] \centering \begin{threeparttable} \caption{Resultados de desempenho do Xorshift} \label{tab:tempos2} \begin{tabular}{|c|c|c|c|c|} \hline \textbf{N bits} & \textbf{Tempo médio (sm)} & \textbf{Tempo mín (sm)} & \textbf{Tempo máx (sm)} & \textbf{Desvio padrão (sm)} \\ \hline 40 & 0.001710 & 0.001400 & 0.004229 & 0.000694 \\ \hline 56 & 0.001480 & 0.001373 & 0.002104 & 0.000171 \\ \hline 80 & 0.001846 & 0.001719 & 0.002383 & 0.000163 \\ \hline 128 & 0.002676 & 0.002164 & 0.005172 & 0.000974 \\ \hline 168 & 0.003205 & 0.003046 & 0.003900 & 0.000206 \\ \hline 224 & 0.003624 & 0.003495 & 0.004421 & 0.000225 \\ \hline 256 & 0.005018 & 0.003886 & 0.017319 & 0.003329 \\ \hline 512 & 0.010519 & 0.008000 & 0.014765 & 0.002494 \\ \hline 1024 & 0.015297 & 0.013844 & 0.022194 & 0.002086 \\ \hline 2048 & 0.029532 & 0.027508 & 0.036664 & 0.002539 \\ \hline 4096 & 0.060360 & 0.058132 & 0.081999 & 0.005983 \\ \hline \end{tabular} \begin{tablenotes} \small \item Os valores representam a média, mínimo, máximo e desvio padrão de 16 execuções independentes para cada tamanho de chave. \end{tablenotes} \end{threeparttable} \end{table}

\subsubsection{Análise dos Resultados}

Em ambos os algoritmos, pode-se perceber que, conforme o número de bits aumenta, ou seja, aumenta-se a complexidade, o desempenho tende a piorar. De forma geral, ambos desempenham similarmente para números menores, aproximadamente até 256 bits, com na casa dos microssegundos e praticamente desprezíveis; entretanto, a diferença é mais notável a partir de 512 bits, no qual o tempo do Xorshift começa a crescer em taxa mais notável que o do LCG, sendo o último o que melhor desempenhou para o caso de 4096 bits.

A diferença na execução dos algoritmos, ainda que com computacional bastante semelhante, possui pontos que tornam o LCG mais performático em arquiteturas de CPUs modernas. Ainda que o Xorshift utilize operações \textit{bitwise}, que individualmente são mais rápidas, acumulam \textit{overhead} por outros pontos do algoritmo, como um estado com variáveis que pode resultar em mais \textit{cache miss}, além de essas operações serem mais difíceis de serem otimizadas em conjunto, bem como possui mais operações por ciclo de cálculo. Enquanto isso, o LCG usa operações de adição, módulo e multiplicação, que se beneficiam de otimizações de CPU, além de ter acesso a memória mais previsível. Somado a isso, o LCG possui complexidade assintótica com mais baixas, tornando no longo prazo o algoritmo mais eficiente para números grandes.

De forma geral, pode-se destacar para cada algoritmo:

\begin{itemize}
    \item \textbf{LCG}
    \begin{itemize}
        \item Pontos positivos:
        \begin{enumerate}
            \item Melhor desempenho com grandes: A implementação se mostrou mais rápida que o Xorshift para números de 512 bits ou mais. Isso ocorre porque a linguagem Python possui suporte nativo e otimizado para aritmética de precisão arbitrária ("bignum"). As operações de multiplicação e adição em números muito grandes são altamente eficientes nessa implementação, o que favorece diretamente o LCG.
            \item Simplicidade de Implementação: O algoritmo é conceitualmente simples e direto de ser implementado.
        \end{enumerate}
        \item Pontos negativos:
        \begin{enumerate}
            \item Qualidade estatística: A qualidade dos valores gerados é extremamente sensível à escolha dos parâmetros iniciais, como tratado na subseção \ref{sub-lcg}. Uma escolha inadequada pode levar a correlações e padrões nos números gerados.
        \end{enumerate}
    \end{itemize}

    \item \textbf{Xorshift} 
    \begin{itemize}
        \item Pontos positivos:
        \begin{enumerate}
            \item Velocidade teórica: O algoritmo é projetado para ser excepcionalmente mais rápido em hardware que executa operações bitwise em poucos ciclos de clock, especialmente números inteiros de tamanho nativo (como 64 bits).
            \item Boas propriedades estatísticas: o algoritmo e suas variantes podem passar por rigorosas baterias de testes estatísticos
        \end{enumerate}

        \item Pontos negativos:
        \begin{enumerate}
            \item Baixo Desempenho com" em Python: O ponto fraco da implementação foi seu desempenho inferior com grandes. A razão é que as operações bitwise (como << ou \^) não são executadas em um único ciclo de máquina em inteiros de 4096 bits. Em Python, essas operações são emuladas sobre uma estrutura de dados que representa o número grande, o que acarreta uma sobrecarga computacional significativa e acaba sendo mais lento do que as operações aritméticas otimizadas do LCG para a mesma escala de número.
        \end{enumerate}
    \end{itemize}
\end{itemize}

\subsection{Avaliação Estatística da Qualidade dos PRNGs}

Para uma avaliação completa da qualidade dos geradores implementados, foi desenvolvida
uma bateria de testes estatísticos específicos para PRNGs. Esta avaliação é crucial
pois, em aplicações criptográficas, a qualidade estatística da sequência
pseudo-aleatória afeta diretamente a segurança do sistema. Os testes implementados
seguem metodologias reconhecidas na literatura e incluem verificações de
reprodutibilidade, distribuição de bits, uniformidade estatística e detecção de
correlações.

\subsubsection{Metodologia dos Testes Estatísticos}

Os testes estatísticos foram organizados em quatro categorias principais:

\begin{enumerate}
    \item \textbf{Testes Básicos}: Verificação de reprodutibilidade (mesma seed produz
mesma sequência) e balanceamento de bits (proporção próxima a 50\% para cada posição de
bit).

    \item \textbf{Testes de Formato}: Validação do método \texttt{randbits(k)} para
diferentes valores de k (1, 4, 8, 16, 32, 64 bits), verificando se os números gerados
possuem o formato correto com bit mais significativo sempre igual a 1.

    \item \textbf{Testes Estatísticos}:
    \begin{itemize}
        \item \textbf{Teste Chi-Quadrado}: Avalia a uniformidade da distribuição
dividindo o espaço [0,1) em 100 intervalos e testando se a frequência observada difere
significativamente da esperada. O valor crítico utilizado foi 124,342 (99 graus de
liberdade, $\alpha$ = 0,05).
        \item \textbf{Teste de Runs}: Converte a sequência em valores binários
(acima/abaixo da mediana) e conta o número de "runs" (sequências consecutivas de 0s ou
1s). Usa estatística z com de confiança de 95\%.
        \item \textbf{Teste de Correlação Serial}: Calcula o coeficiente de correlação
de Pearson entre valores consecutivos da sequência (lag=1) para detectar dependências
lineares.
    \end{itemize}

    \item \textbf{Teste de Período}: Tentativa de detectar repetições no estado interno
do gerador, fundamental para avaliar o período efetivo.
\end{enumerate}

\subsubsection{Resultados dos Testes Estatísticos}

Os experimentos foram conduzidos com fixa (12345) e amostra de 10.000 valores para
os testes estatísticos. Os resultados obtidos estão sumarizados na Tabela
\ref{tab:statistical_tests}.

\begin{table}[h!]
\centering
\begin{threeparttable}
\caption{Resultados dos testes estatísticos para os PRNGs implementados}
\label{tab:statistical_tests}
\begin{tabular}{|l|c|c|}
\hline
\textbf{Teste} & \textbf{LCG} & \textbf{XORShift} \\
\hline
\multicolumn{3}{|c|}{\textit{Testes Básicos}} \\
\hline
Reprodutibilidade & X & X \\
Balanceamento de Bits & X & X \\
\hline
\multicolumn{3}{|c|}{\textit{Testes de Formato (randbits)}} \\
\hline
1, 4, 8, 16, 32, 64 bits & X (todos) & X (todos) \\
\hline
\multicolumn{3}{|c|}{\textit{Testes Estatísticos}} \\
\hline
Chi-Quadrado & X (123,200) & X (93,640) \\
Teste de Runs & X (512 runs) & X (517 runs) \\
Correlação Serial & X (-0,033) & X (-0,017) \\
\hline
\multicolumn{3}{|c|}{\textit{Teste de Período}} \\
\hline
Detecção de Período & Não detectado* & Não detectado* \\
\hline
\multicolumn{3}{|c|}{\textit{Avaliação Geral}} \\
\hline
Taxa de Sucesso & 100\% (11/11) & 100\% (11/11) \\
Status & Excelente & Excelente \\
\hline
\end{tabular}
\begin{tablenotes}
\small
\item * Período não encontrado nas primeiras 100.000 iterações, indicando período longo.
\item Os valores entre parênteses representam as estatísticas calculadas para cada
teste.
\end{tablenotes}
\end{threeparttable}
\end{table}

\subsubsection{Análise dos Resultados Estatísticos}

Ambos os geradores demonstraram excelente desempenho nos testes estatísticos
implementados, aprovando em todos os 11 testes realizados. Esta performance indica que,
do ponto de vista da qualidade estatística básica, tanto o LCG quanto o XORShift são
adequados para a geração de candidatos a números primos.

\textbf{Análise por categoria de teste:}

\begin{itemize}
    \item \textbf{Uniformidade}: O teste chi-quadrado mostrou que ambos os geradores
produzem distribuições estatisticamente uniformes. O XORShift apresentou valor
ligeiramente melhor (93,640) comparado ao LCG (123,200), ambos bem abaixo do valor
crítico de 124,342.

    \item \textbf{Independência}: Os testes de runs e correlação serial confirmaram
ausência de padrões óbvios nas sequências geradas. Os coeficientes de correlação serial
próximos a zero (-0,033 para LCG e -0,017 para XORShift) indicam baixa dependência
linear entre valores consecutivos.

    \item \textbf{Período}: Nenhum dos geradores apresentou repetição de estado nas
primeiras 100.000 iterações, sugerindo períodos longos adequados para as aplicações
pretendidas.
\end{itemize}

É importante ressaltar que, embora estes resultados sejam encorajadores, os testes
implementados representam apenas uma avaliação básica da qualidade estatística. Para
aplicações criptográficas críticas, seria recomendável submeter os geradores a baterias
de testes mais rigorosas, com testes estatísticos mais sofisticados e que
são capazes de detectar deficiências sutis que podem não ser aparentes em análises mais
simples.

\section{Verificação de Primalidade}

Após a geração de um número candidato, a segunda etapa consiste em determinar eficientemente se ele é primo. Dada a magnitude dos números envolvidos, métodos determinísticos como a divisão por tentativa são computacionalmente inviáveis. A solução reside no uso de testes de primalidade probabilísticos, que oferecem um equilíbrio entre velocidade e uma probabilidade de erro controlável e extremamente baixa.

\subsection{Testes de Primalidade de Miller-Rabin}

O Teste de Miller-Rabin é um dos algoritmos de primalidade probabilísticos mais eficientes e amplamente utilizados na prática, sendo um padrão em bibliotecas criptográficas. A versão original, proposta por Gary Miller em 1976, era determinística, mas sua correção dependia da Hipótese de Riemann Generalizada, que permanece não provada. Em 1980, Michael Rabin modificou o algoritmo para uma versão probabilística incondicional, que é a forma utilizada hoje. \parencite{rabin1980}

O teste é uma sofisticação do Teste de Fermat e se baseia em duas propriedades fundamentais dos números primos ímpares n: \parencite{kleinberg2010miller}
\begin{enumerate}
    \item \textbf{Pequeno Teorema de Fermat}: para qualquer inteiro $a$ tal que $1 < a < n-1$, a congruência $a^{n-1} \equiv 1 \;(mod \,n)$ é válida.
    \item \textbf{Raízes quadradas das Unidades}: as únicas soluções para a equação $x^2 \equiv 1 \; (mod \, n)$ são $x \equiv -1 \; (mod \, n)$ e $x \equiv -1 \; (mod \, n)$. Se for encontrada outra raíz quadrada de 1 ("raíz quadrada não trivial"), então n é garantidamente composto. \parencite{wikipedia_miller_rabin}
\end{enumerate}

O algoritmo de Miller-Rabin combina essas duas propriedades de forma engenhosa. Para testar um número ímpar n, primeiro decompõe-se $n-1$ na forma $2^s \cdot d$, onde d é ímpar. Em seguida, escolhe-se um inteiro aleatório a no intervalo $[2, n-2]$, chamado "base" ou "testemunha". O número n passa no teste para a base "a" se uma das condições sejam satisfeitas:
\begin{itemize}
    \item $a^d \equiv1 \; (mod \, n)$
    \item $a^{2^r \cdot d} \equiv -1 \; (mod \, n)$ para algum $r$ no intervalo $0 \leq r < s$.
\end{itemize}

Se n for primo, ele passará no teste para qualquer base a. Se n for composto, ele passará no teste para no máximo 1/4 das bases possíveis. Os números a para os quais um n composto passa no teste são chamados de "mentirosos fortes" (strong liars). Ao repetir o teste k vezes com aleatórias independentes, a probabilidade de que um número composto seja incorretamente classificado como primo é de no máximo $(1/4)^k$. Este erro pode ser tornado arbitrariamente pequeno aumentando o número de iterações k, tornando o teste extremamente confiável para aplicações criptográficas. 

\subsection{Teste de Primalidade de Fermat}

A escolha do Teste de Fermat como segundo método para comparação é deliberada. Sua simplicidade conceitual e, mais importante, sua falha teórica bem conhecida, servem como um excelente contraponto para ilustrar a robustez e a superioridade do Teste de Miller-Rabin.


O teste baseia-se diretamente e unicamente no Pequeno Teorema de Fermat. Para testar um número n, escolhe-se um uma base aleatória $a$ (com < a < n-1) e verifica-se a congruência $a^{n-1} \equiv 1 \; (mod \, n)$. \parencite{ma2014fermat}

\begin{itemize}
    \item Se $a^{n-1} \not\equiv 1 \; (mod \, n)$, então a é uma "testemunha de Fermat" para a composição de n, e n é definitivamente composto.
    \item Se $a^{n-1} \equiv 1 \; (mod \, n)$, então n é "provavelmente primo" para a base a e a é uma "mentirosa de Fermat".
\end{itemize}

principal fraqueza do Teste de Fermat é a existência de números compostos que são mentirosos para todas as bases a que são coprimas com. Esses números são chamados de Números de Carmichael. O menor deles é 561 = $3 \cdot 11 \cdot 17$. Para qualquer a coprimo com, a congruência $a^{560} \equiv 1 \; (mod \, 561)$ é satisfeita, fazendo com o teste falhe em detectar a composição de 561. \parencite{wikipedia_pequeno_teorema_fermat}. Como existem infinitos números de Carmichael, o Teste de Fermat não pode ser usado como um teste de primalidade de propósito geral confiável. Em contraste, o Teste de Miller-Rabin foi projetado especificamente para detectar a composição de números como os de Carmichael, explorando a propriedade das raízes quadradas não triviais da unidade.

\subsection{Implementação e Metodologia Experimental}
A implementação e experimentação dos testes de primalidade foram projetadas para avaliar os algoritmos no contexto de números inteiros de grande porte, conforme exigido pelas aplicações criptográficas modernas.

\subsubsection{Código Fonte e Metodologia}

Os experimentos foram conduzidos com mesmo hardware e demais especificações da seção \ref{sec:config-exp}. Os testes foram também implementados em Python, com fonte disponível no GitHub (Apêndice \ref{app-a}). A metodologia experimental consistiu em:
\begin{itemize}
    \item Para cada tamanho de bits (40, 56, 80, 128, 168, 224, 256, 512, 1024, 2048, 4096), utilizou-se o PRNG LCG para gerar números candidatos.
    \item Cada candidato foi testado por ambos algoritmos de primalidade, Miller-Rabin com=20 iterações e Fermat com=20 bases aleatórias.
    \item O processo continua, gerando e testando novos candidates até que o primeiro número primo seja encontrado por um dos testes.
    \item O tempo total decorrido para encontrar o primeiro primo, bem como o número de candidatos testados, foram registrados.
\end{itemize}

\subsubsection{Experimentos}

Os resultados obtidos estão sumarizados na Tabela \ref{tab:prime_discovery}. A tabela apresennta o tempo total decorrido para encontrar o primeiro número primo, o número de candidatos testados, e se cada teste (Miller-Rabin, Fermat ou ambos) conseguiu identificar corretamente o primo.

\begin{table}[h!]
    \centering
    \begin{threeparttable}
    \caption{Resultados do experimento de descoberta de números primos}
    \label{tab:prime_discovery}
    \begin{tabular}{|c|c|c|c|c|c|}
    \hline
    \textbf{N bits} & \textbf{Tempo (ms)} & \textbf{Candidatos} & \textbf{Miller-Rabin} &
    \textbf{Fermat} & \textbf{Ambos} \\
    \hline
    40 & 0.234154 & 3 & Sim & Sim & Sim \\
    \hline
    56 & 1.734929 & 103 & Sim & Sim & Sim \\
    \hline
    80 & 0.863268 & 17 & Sim & Sim & Sim \\
    \hline
    128 & 10.825385 & 160 & Sim & Sim & Sim \\
    \hline
    168 & 5.039106 & 37 & Sim & Sim & Sim \\
    \hline
    224 & 7.976166 & 31 & Sim & Sim & Sim \\
    \hline
    256 & 29.520978 & 130 & Sim & Sim & Sim \\
    \hline
    512 & 39.746970 & 18 & Sim & Sim & Sim \\
    \hline
    1024 & 485.609992 & 60 & Sim & Sim & Sim \\
    \hline
    2048 & 12817.833151 & 322 & Sim & Sim & Sim \\
    \hline
    4096 & 76543.884109 & 276 & Sim & Sim & Sim \\
    \hline
    \end{tabular}
    \begin{tablenotes}
    \small
    \item Tempo decorrido para encontrar o primeiro número primo usando LCG para geração de
    candidatos ímpares e testes de Miller-Rabin e Fermat (k=20). Seed fixo: 42.
    \end{tablenotes}
    \end{threeparttable}
    \end{table}
  

\subsubsection{Análise dos Resultados}
Os resultados indicam que ambos os testes de primalidade foram eficazes em identificar números primos para todos os tamanhos de bits testados. No entanto, o tempo necessário para encontrar o primeiro primo aumentou significativamente com tamanho dos números, refletindo a complexidade crescente do problema. Como esperado, os tempos de execução para encontrar um primo são muito semelhantes para ambos os algoritmos. Isso ocorre porque o custo computacional de uma única iteração de cada teste é dominado pela operação de exponenciação modular, que é idêntica em ambos os casos. A complexidade teórica se reflete nos dados empíricos, onde o tempo de execução aumenta drasticamente com número de bits. A dificuldade crescente em encontrar primos maiores também se deve à sua menor densidade, conforme previsto pelo Teorema dos Números Primos, que afirma que a probabilidade de um número aleatório N ser primo é aproximadamente $\frac{1}{\ln(N)}$. \parencite{wikipedia_numero_primo}

\subsubsection{Análise de Robustez e Correção}

A análise de desempenho por si só é enganosa, pois não revela a diferença fundamental entre os dois testes. A verdadeira distinção reside na sua robustez. Para demonstrar isso empiricamente, foi realizado um experimento direcionado, testando ambos os algoritmos com números conhecidos por serem casos problemáticos para o Teste de Fermat.

O enunciado do trabalho sugere o uso de pseudoprimos. A escolha do Teste de Fermat permite um teste ainda mais contundente: o uso de Números de Carmichael. Fontes como a \textit{On-Line Encyclopedia of Integer Sequences} (OEIS) e o MathWorld fornecem listas desses números. Os primeiros são 561, 1105 e 1729.

Para demonstrar adequadamente a vulnerabilidade do Teste de Fermat, foi implementado um experimento específico que testa os números de Carmichael usando apenas bases coprimas com o número candidato. Este cenário representa o caso mais desfavorável para o Teste de Fermat, onde sua vulnerabilidade teórica se manifesta na prática.

Os resultados obtidos com o experimento implementado em \texttt{carmichael\_test.py} demonstram inequivocamente a falha teórica do Teste de Fermat:

\begin{itemize}
    \item Para o número 561 = $3 \times 11 \times 17$:
    \begin{itemize}
        \item \texttt{fermat\_test} (bases aleatórias, k=20): \textbf{False} (detecta corretamente como composto)
        \item \texttt{fermat\_test} (apenas bases coprimas, k=20): \textbf{True} (falha - classifica como primo)
        \item \texttt{miller\_rabin\_test} (k=20): \textbf{False} (detecta corretamente como composto)
    \end{itemize}

    \item Resultados similares foram obtidos para 1105 = $5 \times 13 \times 17$ e 1729 = $7 \times 13 \times 19$.
\end{itemize}

Este resultado prático é uma demonstração inequívoca da falha teórica do Teste de Fermat. Quando testado apenas com bases coprimas, ele classifica incorretamente os números compostos de Carmichael como "provavelmente primo" porque, para essas bases, a congruência $a^{n-1} \equiv 1 \pmod{n}$ é sempre satisfeita. Por outro lado, o Teste de Miller-Rabin identifica corretamente todos os números de Carmichael como compostos.

É importante notar que na implementação padrão do Teste de Fermat com bases aleatórias, a probabilidade de escolher uma base coprima diminui conforme o número tem mais fatores pequenos. Para 561, as bases não coprimas incluem todos os múltiplos de 3, 11 e 17, representando uma fração significativa do espaço de busca. Isso explica por que o teste com bases aleatórias frequentemente detecta 561 como composto, mas não elimina a vulnerabilidade teórica fundamental.

Esta vulnerabilidade torna o Teste de Fermat inadequado para aplicações criptográficas que exigem alta confiança. O Teste de Miller-Rabin, com sua probabilidade de erro garantida de no máximo $(1/4)^k$, oferece o nível de segurança necessário. Para $k=20$ iterações, a probabilidade de um número composto ser classificado como primo é inferior a 1 em $10^{12}$, um nível de confiança mais do que suficiente para a maioria das aplicações práticas.


\section{Conclusão}

Este trabalho realizou uma análise teórica e empírica de algoritmos para a geração de números pseudo-aleatórios e para a verificação de primalidade, duas etapas fundamentais na construção de sistemas criptográficos modernos. A implementação e a avaliação do Gerador Linear Congruencial (LCG), do Xorshift, do Teste de Miller-Rabin e do Teste de Fermat com inteiros de até 4096 bits permitiram extrair conclusões significativas sobre o compromisso entre desempenho, complexidade teórica e robustez prática.

Na geração de números pseudo-aleatórios, os resultados experimentais revelaram uma inversão de desempenho notável em relação às expectativas teóricas. Embora o Xorshift seja projetado para ser mais rápido por utilizar operações bitwise de baixo custo, o LCG demonstrou ser mais performático para números de grande magnitude no ambiente Python. Esta discrepância foi atribuída à implementação nativa altamente otimizada da aritmética de precisão arbitrária ("bignum") do Python, que beneficia diretamente as operações de multiplicação e módulo do LCG, em contraste com a sobrecarga computacional da emulação de operações bitwise em inteiros tão grandes. No que tange à qualidade, ambos os geradores foram aprovados com sucesso em uma bateria de testes estatísticos básicos, mostrando-se adequados para a geração de números candidatos. 

Na verificação de primalidade, a análise comparativa entre o Teste de Miller-Rabin e o Teste de Fermat confirmou que o desempenho de ambos é praticamente idêntico, visto que o custo computacional de ambos é dominado pela operação de exponenciação modular.  Contudo, a diferença crucial e fundamental reside na robustez. O experimento direcionado com números de Carmichael demonstrou inequivocamente a vulnerabilidade teórica do Teste de Fermat, que classificou incorretamente um número composto como primo sob condições específicas. Em contrapartida, o Teste de Miller-Rabin provou sua superioridade ao identificar corretamente a composição em todos os cenários.

Em suma, o estudo evidencia que, enquanto a performance de algoritmos pode variar de acordo com o ambiente de implementação de maneiras não triviais, a robustez teórica é um critério inegociável para aplicações de segurança. O Teste de Miller-Rabin, com sua confiabilidade matemática garantida, reafirma-se como a escolha indispensável para a criptografia, ao passo que a seleção de um PRNG deve ponderar a velocidade teórica do algoritmo frente às otimizações da plataforma de implementação.


% --- BIBLIOGRAFIA --- 

\newpage

\printbibliography

\newpage
\appendix
\section{Apêndice A - Código Fonte}
\label{app-a}

O código fonte para o trabalho, assim como os experimentos realizados e testes, pode ser encontrado no repositório público no GitHub, disponível em: \url{https://github.com/caioopra/prime-numbers}

\end{document}

